\section{An authority contract}
\label{sec:contract}

\subsection{Scope before probability}
An authority contract records a deployment's permitted effects and the evidence supporting them. It is an engineering record, not a legal contract or a claim that an AI has agreed to anything. At minimum it specifies:
\begin{enumerate}
\item the task, affected parties, protected interests, and forbidden outcomes;
\item the grounds and scope of each approving party\textquotesingle s authority, and the process for challenge;
\item the model, tools, runtime, memory policy, and environment versions;
\item the action interface, actors, credentials, and human-mediated channels;
\item resource, concurrency, duration, and cumulative-exposure limits;
\item enforced invariants, empirical findings, and unresolved assumptions;
\item revocation, recovery, independent ownership, and revalidation triggers.
\end{enumerate}
Both misuse by people and deviation by the model belong in the threat model. A system that obeys an unauthorized harmful request is not made safe by that obedience. The contract must encode legitimate scope separately from what any one requester wants.


\subsection{Authorization is specific to an effect}
For a proposed effect $e$ in state $x$, let $J(e,x)$ identify the authorizations required by the governing specification. Let $g_j(e,x)$ state whether issuer $j$ has a valid, in-scope grant for that effect. A basic execution condition is
\begin{equation}
\operatorname{Permit}(e,x)=C(e,x)\wedge\bigwedge_{j\in J(e,x)}g_j(e,x),
\label{eq:jointgrant}
\end{equation}
where $C$ includes resource, version, and other execution constraints. Required grants compose by conjunction, not by a vote over interchangeable reviewers. The default for an unresolved required authorization is to withhold that effect. This rule does not imply that every affected party is a required issuer: determining $J$ is a prior institutional obligation, not something the proposer may choose.

For example, permission to spend compute does not confer permission to reuse a community-controlled dataset. A valid operator grant and a missing required data-use grant must yield no execution. A change of purpose must trigger a fresh scope check. Conversely, copying an actor must not duplicate the authority or budget of its issuer. The finite model in Section~\ref{sec:experiment} has one outside controller. The SQLite prototype additionally requires exact-payload operator and steward approvals for synthetic data release and document edits. This checks conjunction of encoded grants; it does not test multi-party governance or issuer legitimacy. The invariant below can enforce a correctly encoded rule without proving that the chosen issuers or their decisions are legitimate.

\subsection{A policy-independent invariant}
Let $x_t$ be the modeled state, including authorization, cumulative resource use, pending operations, and relevant environmental state. An unrestricted proposer emits $a_t\in\A$. An external enforcer chooses an executed operation or safe fallback $E(x_t,a_t)$. Disturbances lie in a specified set $W$, and
\[
x_{t+1}\in F(x_t,E(x_t,a_t),w_t),\qquad w_t\in W.
\]
Let $\I$ be an invariant and $\Hh$ the modeled harmful states, with $\I\cap\Hh=\varnothing$.

\begin{lemma}[Boundary invariant; standard induction]
\label{lem:invariant}
If $x_0\in\I$ and
\[
F(x,E(x,a),w)\subseteq\I
\quad\text{for all }x\in\I,\ a\in\A,\ w\in W,
\]
then every modeled execution remains in $\I$, whatever policy chooses the proposals.
\end{lemma}
\begin{proof}
The initial state satisfies the invariant. The displayed inclusion preserves it at each transition. Induction gives the result for every finite prefix, hence every state of any infinite execution.
\end{proof}

This standard construction is useful because the proposer does not have to be trustworthy. The universal quantifier already includes adaptive, deceptive, mistaken, or very capable proposal policies. It does not include transitions omitted from $F$.

\para{From a model to a deployment.} Within an institutionally justified scope, three technical obligations carry the practical argument: \emph{specification adequacy} (the property addresses the intended harm); \emph{refinement} (the implementation's effects correspond to the modeled transitions); and \emph{complete mediation} (there is no accessible route around that implementation). The trusted base includes the relevant operating system, isolation, compiler, hardware assumptions, and authorization owners. Formal verification can reduce this base; it cannot erase it by definition. Verified kernels demonstrate substantial progress on such bounded obligations \cite{sel4}.

\para{Capability and authority.} Define the modeled authority of an interface as the reachable effects under \emph{all} proposal sequences, from the stated initial states. Holding dynamics and initial states fixed, deleting allowed operations can only shrink that set: every remaining trace was already possible before the deletion. This monotonicity is a structural result. It is not a claim that removing one permission always reduces measured harm in a changing real organization; operators may compensate through other channels.

An upgrade can change the operational threat even when the API stays identical. More copies, faster requests, longer memory, or stronger persuasion may enlarge effective exposure. Therefore ``no automatic authority expansion'' must cover aggregate resources and channels, not merely unchanged permission names.



\subsection{What survives an upgrade}
\label{sec:upgrade}
Because Lemma~\ref{lem:invariant} already quantifies over all proposals, replacing one policy with another, or switching policies using the execution history, does not change its conclusion. This is an immediate corollary, not a theorem about arbitrary software upgrades. It requires the same transition model, enforced invariant, initial-state conditions, and modeled disturbance set. More parallel actors must already be represented in that model; a new tool, credential, communication channel, or execution path creates a new obligation.

Evidence therefore has a validity domain, not just a date. A policy upgrade may preserve a checked budget invariant while invalidating a monitor's measured failure interval. A runtime upgrade may preserve task accuracy while invalidating refinement evidence. Renewal must name which premises changed and discharge the affected obligations. Re-running one aggregate benchmark cannot stand in for that analysis. Table~\ref{tab:renewal} makes the distinction operational.

\begin{table*}[t]
\centering\small
\begin{tabularx}{\textwidth}{@{}P{.19\textwidth}LL@{}}
\toprule
\textbf{Change} & \textbf{Evidence that may remain valid} & \textbf{Obligation before authority continues}\\
\midrule
Proposer policy or weights & An all-policy invariant under the same modeled interface and dynamics & Reassess behavioral and monitor evidence; confirm that tools, channels, and implementation did not change\\
Copies, speed, or duration & Invariants whose actor and scheduling model already covers the change & Preserve the aggregate ledger; recompute exposure and contention or timing assumptions\\
Runtime, tools, or credentials & Unchanged abstract properties and their proofs & Re-establish implementation refinement and complete mediation for the new effects\\
Restart or renewed grant & Evidence whose premises persist through the transition & Preserve consumption and identifier history; recheck liveness and version at commit\\
Protected interest or specification & Evidence for unaffected properties only & Validate the new predicates independently; replay affected checks and useful-task witnesses\\
New human-mediated channel & Only claims that already covered that channel & Assess induced actions and decision authority; retain unknown scope where mediation is unproved\\
\bottomrule
\end{tabularx}
\caption{Renew evidence at the level of its premises. Retaining a valid result and withholding an invalid extrapolation are compatible. This is a proposed decision protocol, not an observed deployment outcome.}
\label{tab:renewal}
\end{table*}

\begin{lemma}[Persistent allowance]
\label{lem:ledger}
Let $z$ be cumulative executed consumption and $L$ the ledger used by an enforcer, initially both zero. An atomic commit of $d\ge0$ requires $L+d\le B$ and increments both by $d$. Suppose all other modeled transitions, including restart and renewal, preserve both. Then $z=L\le B$ for every execution, irrespective of the number of actors, proposals, or renewals.
\end{lemma}
\begin{proof}
The equality and inequality hold initially. A commit preserves equality and checks the new bound. Every other transition preserves both quantities. Induction completes the argument.
\end{proof}
A genuine reversal or increase in authorized budget needs an explicit transition and a new obligation; it cannot be implemented by erasing the ledger. \emph{Renewal restores permission; it does not refund effects.} The experiment below distinguishes real cumulative consumption from the enforcer's accounting precisely to test this failure.

\subsection{Safety must leave room for useful work}
Rejecting every request preserves many invariants. It is often operationally worthless and may push people toward uncontrolled alternatives. An assurance argument therefore needs positive witnesses: legitimate tasks that remain executable, with their quality, cost, and delay. These witnesses test non-vacuity; they do not prove that every important user need is supported.

The design objective is a constrained usefulness frontier: improve useful outcomes while respecting the stated invariant and exposure budget. Training remains valuable for selecting good proposals, reducing unnecessary denials, and improving the evidence presented to reviewers. Nothing in the boundary argument makes aligned behavior irrelevant to usefulness or to risks that the boundary cannot express.
